\par\Needspace{6\baselineskip}\medskip
\hypertarget{native-091-article071}{}\pdfbookmark[2]{71.}{article-071}
\begin{center}71.\end{center}
Aberratio fixarum, nec non pars ea aberrationis planetarum et cometarum, quae soli motui terrae debetur, oritur inde, quod cum terra integra \emph{tubus} movetur, dum radius luminis ipsius axem opticum percurrit. Corporis coelestis locus observatus (qui et apparens seu aberratione affectus dicitur) determinatur per situm axis optici telescopii ita collocati, ut radius luminis ab illo egressus in via sua utramque huius axis extremitatem attingat: hic autem situs diversus est a situ vero radii luminis in spatio. Distinguamus duo temporis momenta $t$, $t'$, ubi radius luminis extremitatem anteriorem (centrum vitri obiectivi), ubique posteriorem (focum vitri obiectivi) attingit; sint harum extremitatum loci in spatio pro momento priori $a$, $b$; pro posteriori $a'$, $b'$. Tunc patet, rectam $ab'$ esse situm verum radii in spatio, loco apparenti autem respondere rectam $ab$ vel $a'b'$ (quas pro parallelis habere licet): nullo porro negotio \SourceLeaf{092}perspicitur, locum apparentem a longitudine tubi non pendere. Differentia inter situm rectarum $b'a$, $ba$ est aberratio qualis pro stellis fixis locum habet: modum eam calculandi hic tamquam notum silentio transimus. Pro stellis errantibus autem ista differentia nondum est aberratio completa: planeta scilicet, dum radius ex ipso egressus ad terram descendit, locum suum ipse mutat, quapropter situs huius radii non respondet loco geocentrico vero tempore observationis. Supponamus, radium luminis qui tempore $t$ in tubum impingit tempore $T$ e planeta egressum esse; designeturque locus planetae in spatio tempore $T$ per $P$, tempore $t$ autem per $p$; denique sit $A$ locus extremitatis antecedentis axis tubi pro tempore $T$. Tunc patet
\par $1^{\mathrm{o}}$ rectam $AP$ exhibere locum verum planetae tempore $T$;
\par $2^{\mathrm{o}}$ rectam $ap$ autem locum verum tempore $t$;
\par $3^{\mathrm{o}}$ rectam $ba$ vel $b'a'$ locum apparentem tempore $t$ vel $t'$ (quorum differentia ceu quantitas infinite parva spectari potest);
\par $4^{\mathrm{o}}$ rectam $b'a$ eundem locum apparentem ab aberratione fixarum purgatum.

Iam puncta $P$, $a$, $b'$ in linea recta iacent, eruntque partes $Pa$, $ab'$ proportionales temporum intervallis $t-T$, $t'-t$, siquidem motus luminis celeritate uniformi peragitur. Temporis intervallum $t'-T$ propter immensam luminis velocitatem semper est perparvum, intra quod motum terrae tamquam rectilineum ac celeritate uniformi peractum supponere licet: sic etiam $A$, $a$, $a'$ in directum iacebunt, partesque $Aa$, $aa'$ quoque intervallis $t-T$, $t'-t$ proportionales erunt. Hinc facile concluditur, rectas $AP$, $b'a'$ esse parallelas, adeoque locum primum cum tertio identicum.

Tempus $t-T$ erit productum distantiae $Pa$ in $493^{\mathrm{s}}$, intra quod lumen percurrit distantiam mediam terrae a Sole, quam pro unitate accepimus. In hoc calculo pro distantia $Pa$ etiam $PA$ vel $pa$ accipere licebit, quum differentia nullius momenti esse possit.

Ex his principiis tres demanant methodi, planetae vel cometae locum apparentem pro quovis tempore $t$ determinandi, e quibus modo hanc modo illam praeferre conveniet.

I. Subtrahatur a tempore proposito tempus intra quod lumen a planeta ad terram descendit: sic prodibit tempus reductum $T$, pro quo locus verus more solito computatus cum apparente pro $t$ identicus erit. Ad computum \SourceLeaf{093}reductionis temporis $t-T$ distantiam a terra novisse oportet: plerumque ad hunc finem subsidia commoda non deerunt, e. g. per ephemeridem vel levi tantum calamo calculatam, alioquin distantiam veram pro tempore $t$ more solito sed neglecta praecisione nimia per calculum praeliminarem determinare sufficiet.

II. Computetur pro tempore proposito $t$ locus verus atque distantia, ex hac reductio temporis $t-T$, atque hinc adiumento motus diurni (in longitudine et latitudine vel in ascensione recta et declinatione) reductio loci veri ad tempus $T$.

III. Computetur locus heliocentricus terrae quidem pro tempore $t$: locus heliocentricus planetae autem pro tempore $T$: dein ex horum combinatione more solito locus geocentricus planetae, qui aberratione fixarum (per methodum notam eruenda sive e tabulis depromenda) auctus locum apparentem quaesitum suppeditabit.

Methodus secunda, quae vulgo in usum vocari solet, eo quidem prae reliquis se commendat, quod ad distantiam determinandam nunquam opus est calculo duplici, attamen eo laborat incommodo, quod adhiberi nequit, nisi plures loci vicini vel calculentur vel ex observationibus iam innotuerint; alioquin enim motum diurnum pro dato habere non liceret.

Incommodum, quo methodus prima et tertia premuntur, plane tollitur, quoties plures loci sibi vicini calculandi sunt. Quam primum enim pro quibusdam distantiae iam innotuerunt, percommode et praecisione sufficiente distantias proxime sequentes per subsidia trita concludere licebit. Ceterum si distantia est nota, methodus prima tertiae ideo plerumque praeferenda erit, quod aberratione fixarum opus non habet; sin vero ad calculum duplicem refugiendum est, tertia eo se commendat, quod in calculo altero locus terrae saltem retinendus est.

Sponte iam se offerunt, quae ad problema inversum requiruntur, puta si e loco apparente verus derivandus est. Scilicet secundum methodum I retinebis locum ipsum immutatum, sed tempus $t$, cui locus propositus ut apparens respondet, convertes in reductum $T$, cui idem tamquam verus respondebit. Secundum methodum II retinebis tempus $t$, sed loco proposito adiicies motum intra tempus $t-T$, quasi istum ad tempus $t+(t-T)$ reducere velles. Secundum methodum III locum propositum ab aberratione fixarum liberatum \SourceLeaf{094}tamquam locum verum pro tempore $T$ considerabis, sed terrae locus verus tempori $t$ respondens retinendus est ac si ad istud pertineret. Utilitas methodi tertiae in Libro secundo clarius elucebit.

Ceterum, ne quid desit, adhuc observamus, locum Solis ab aberratione perinde affici ac locum planetae: sed quoniam tum distantia a terra tum motus diurnus propemodum sunt constantes, aberratio ipsa semper valorem tantum non constantem obtinet motui medio Solis in $493^{\mathrm{s}}$ aequalem, adeoque $=20''{,}25$, quae quantitas a longitudine vera subtrahenda est, ut apparens prodeat. Valor aberrationis exactus est in ratione composita distantiae et motus diurni, sive quod eodem redit in ratione inversa distantiae, unde ille valor medius in apogeo $0''{,}34$ diminuendus in perigeo tantumdem augendus esset. Ceterum tabulae nostrae solares aberrationem constantem $-20''{,}25$ iam includunt; quapropter ad obtinendum longitudinem veram tabulari $20''{,}25$ addere oportebit.

